\documentclass[10pt,a4paper]{amsart}
\usepackage{amsmath,amssymb,amsthm,graphicx}
\usepackage{verbatim}
\newenvironment{reference}{\comment}{\endcomment}
%\newenvironment{reference}{}{}
\newenvironment{slogan}{\comment}{\endcomment}
\newenvironment{history}{\comment}{\endcomment}

% For commutative diagrams we use Xy-pic
\usepackage[all]{xy}

% We use 2cell for 2-commutative diagrams.
\xyoption{2cell}
% Theorem environments.
%
\theoremstyle{plain}
\newtheorem{theorem}[subsection]{Theorem}
\newtheorem{proposition}[subsection]{Proposition}
\newtheorem{lemma}[subsection]{Lemma}

\theoremstyle{definition}
\newtheorem{definition}[subsection]{Definition}
\newtheorem{example}[subsection]{Example}
\newtheorem{exercise}[subsection]{Exercise}
\newtheorem{situation}[subsection]{Situation}

\theoremstyle{remark}
\newtheorem{remark}[subsection]{Remark}
\newtheorem{remarks}[subsection]{Remarks}

\numberwithin{equation}{subsection}

% Macros
%
\def\lim{\mathop{\mathrm{lim}}\nolimits}
\def\colim{\mathop{\mathrm{colim}}\nolimits}
\def\Spec{\mathop{\mathrm{Spec}}}
\def\Hom{\mathop{\mathrm{Hom}}\nolimits}
\def\Ext{\mathop{\mathrm{Ext}}\nolimits}
\def\SheafHom{\mathop{\mathcal{H}\!\mathit{om}}\nolimits}
\def\SheafExt{\mathop{\mathcal{E}\!\mathit{xt}}\nolimits}
\def\Sch{\mathit{Sch}}
\def\Mor{\mathop{\mathrm{Mor}}\nolimits}
\def\Ob{\mathop{\mathrm{Ob}}\nolimits}
\def\Sh{\mathop{\mathit{Sh}}\nolimits}
\def\NL{\mathop{N\!L}\nolimits}
\def\CH{\mathop{\mathrm{CH}}\nolimits}
\def\proetale{{pro\text{-}\acute{e}tale}}
\def\etale{{\acute{e}tale}}
\def\QCoh{\mathit{QCoh}}
\def\Ker{\mathop{\mathrm{Ker}}}
\def\Im{\mathop{\mathrm{Im}}}
\def\Coker{\mathop{\mathrm{Coker}}}
\def\Coim{\mathop{\mathrm{Coim}}}

% Boxtimes
%
\DeclareMathSymbol{\boxtimes}{\mathbin}{AMSa}{"02}

%
% Macros for moduli stacks/spaces
%
\def\QCohstack{\mathcal{QC}\!\mathit{oh}}
\def\Cohstack{\mathcal{C}\!\mathit{oh}}
\def\Spacesstack{\mathcal{S}\!\mathit{paces}}
\def\Quotfunctor{\mathrm{Quot}}
\def\Hilbfunctor{\mathrm{Hilb}}
\def\Curvesstack{\mathcal{C}\!\mathit{urves}}
\def\Polarizedstack{\mathcal{P}\!\mathit{olarized}}
\def\Complexesstack{\mathcal{C}\!\mathit{omplexes}}
% \Pic is the operator that assigns to X its picard group, usage \Pic(X)
% \Picardstack_{X/B} denotes the Picard stack of X over B
% \Picardfunctor_{X/B} denotes the Picard functor of X over B
\def\Pic{\mathop{\mathrm{Pic}}\nolimits}
\def\Picardstack{\mathcal{P}\!\mathit{ic}}
\def\Picardfunctor{\mathrm{Pic}}
\def\Deformationcategory{\mathcal{D}\!\mathit{ef}}
\title{Existence and uniqueness of the tensor product of two modules}
\author{The Stacks Project Authors}
\date{stacks-project:algebra.tex:local-snapshot-20261012; GNU FDL 1.2 or later}
\begin{document}
\maketitle
\noindent\textit{Editorial note.} All delivered mathematical text is the author's own, copied from the pinned Stacks Project chapter file algebra.tex; no statement or proof step was rewritten, reordered or removed.\par\medskip
\noindent\textit{Editorial note.} Preamble provenance, as required for this source: the frozen snapshot stages the chapter's own preamble file beside it (preamble.tex, 217 lines, 8116 bytes, the same shared Stacks-wide preamble the published chow and etale-cohomology cards use); the first line of the pinned chapter file is \textbackslash{}input\{preamble\}, so no preamble text lives inside the chapter itself. The card takes that file's lines 8-18 (the verbatim package and the three \textbackslash{}comment-based environments reference/slogan/history at lines 9-12, section 13.2, together with the Xy-pic setup) and its lines 154 to the runtime end of the file (the native theorem declarations and mathematical macros). Not one character was changed, no package was added, no macro was defined, and the excerpt serves only as preamble, never as a span coordinate.\par\medskip
\noindent\textit{Editorial note.} The same-chapter pointers are located at their \textbackslash{}label lines inside the pinned chapter file; the cross-chapter pointers are located in the project's own tag list staged in the snapshot (sources/stacks-chow/tags.txt), at the verbatim lines carrying their canonical tags.\par\medskip
\noindent\textit{Editorial note.} The counted statement contains an Xy-pic diagram; the preamble spans of this card run 8-18 of the chapter's own preamble file and therefore include the Xy-pic package and option setup at lines 15-18.\par\medskip
\noindent\textit{Editorial note.} The counted spans contain no \textbackslash{}ref and no \textbackslash{}cite, so the card carries no navigation stub and no bibliography.\par\medskip
\noindent\textit{Allowed standard prerequisites (editorial).}\begin{itemize}
\item Free module on the product set and its quotient. Used in the counted proof at lines 1963-1965. Supplies the candidate T as a quotient killing the bilinearity relations.
\end{itemize}\medskip
\begin{lemma}
\label{lemma-tensor-product}
Let $M, N$ be $R$-modules. Then there exists a pair $(T, g)$
where $T$ is an $R$-module, and
$g : M \times N \to T$ an $R$-bilinear
mapping, with the following universal property:
For any $R$-module $P$ and any $R$-bilinear mapping
$f : M \times N \to P$, there
exists a unique $R$-linear
mapping $\tilde{f} : T \to P$ such that $f = \tilde{f} \circ g$.
In other words, the following diagram commutes:
$$
\xymatrix{
M \times N \ar[rr]^f \ar[dr]_g & & P\\
& T \ar[ur]_{\tilde f}
}
$$
Moreover, if $(T, g)$ and $(T', g')$
are two pairs with this property, then there
exists a unique isomorphism
$j : T \to T'$ such that $j\circ g = g'$.
\end{lemma}

\begin{proof}
We first prove the existence of such $R$-module $T$.
Let $M, N$ be $R$-modules.
Let $T$ be the quotient module
$P/Q$, where $P$ is the free $R$-module $R^{(M \times N)}$ and $Q$ is the
$R$-module generated by all elements of
the following types: ($x\in M, y\in N$)
\begin{align*}
(x + x', y) - (x, y) - (x', y), \\
(x, y + y') - (x, y) - (x, y'), \\
(ax, y) - a(x, y), \\
(x, ay) - a(x, y)
\end{align*}
Let $\pi : M \times N \to T$ denote the natural map.
This map is $R$-bilinear, as
implied by the above relations
when we check the bilinearity conditions. Denote the image
$\pi(x, y) = x \otimes
y$, then these elements generate
$T$. Now let $f : M \times N \to P$ be an $R$-bilinear map,
then we can define
$f' : T \to P$ by extending the mapping
$f'(x \otimes y) = f(x, y)$. Clearly $f = f'\circ \pi$. Moreover, $f'$ is
uniquely determined by the value on the
generating sets $\{x \otimes y : x\in M, y\in N\}$.
Suppose there is another pair $(T', g')$ satisfying the same properties.
Then there is a unique $j : T \to T'$ and
also $j' : T' \to T$ such that $g' = j\circ g$, $g = j'\circ g'$.
But then both the maps $(j\circ j') \circ g$ and $g$
satisfies the universal properties, so by uniqueness they are equal,
and hence $j'\circ j$ is identity on $T$.
Similarly $(j'\circ j) \circ g' = g'$ and $j\circ j'$ is identity on $T'$.
So $j$ is an isomorphism.
\end{proof}

\end{document}
